% !PRINTED 270
\weberpage{270}
\begin{center}
\gesperrt{Eighth Section.}

\gesperrt{Sturm's Theorem.}

\rule{0.3\linewidth}{0.4pt}

\S.\,85.

\gesperrt{Sturm's Problem.}
\end{center}

The question treated in the previous section, of the number of real roots of an equation, is a special case of a more general problem, which is to form the subject of this section, and which is the foundation of all methods for the approximate numerical computation of the roots of equations.

The question is this: \gesperrt{how many real roots of a real numerical equation lie between two given real values} $a$ and $b$?

Once this is decided, the further task is to narrow the bounds $a$, $b$ until only \gesperrt{one root} of the given equation lies between them, and finally to bring them so close to one another that each of them can be regarded as an \gesperrt{approximate value} of this root.

If we take $a = -\infty$, $b = +\infty$, or at any rate $a$ negative and $b$ positive and so large that no roots can lie beyond these bounds, then this problem coincides with the one treated in the previous section, of determining the number of all the real roots.

We shall first show how the general problem can be reduced to the special one, and thus how the question we have now raised is, in principle at least, answered by the considerations of the previous section. Let us first consider
%CONT>>
